% Glossary entries for the belief-behavior alignment paper
% Add \usepackage{glossaries} to your preamble and \makeglossaries before \begin{document}
% Use \printglossaries where you want the glossary to appear

% Core concepts
\newglossaryentry{belief-behavior-alignment}{
    name={belief-behavior alignment},
    description={The consistency between an agent's explicitly stated beliefs about human behavior and their actual simulated actions in role-playing scenarios}
}

\newglossaryentry{role-playing-agent}{
    name={role-playing agent},
    description={Large language model configured to simulate human behavior by adopting specific personas or demographic characteristics}
}

\newglossaryentry{synthetic-data-generation}{
    name={synthetic data generation},
    description={The process of creating artificial datasets that mimic real human behavioral patterns using computational models}
}

\newglossaryentry{self-consistency}{
    name={self-consistency},
    description={Internal coherence between different outputs or behaviors of the same model, particularly between stated beliefs and enacted behaviors}
}

% Game theory and experimental setup
\newglossaryentry{trust-game}{
    name={Trust Game},
    description={A two-player economic game where Player A (Trustor) sends money to Player B (Trustee), the amount is multiplied, and Player B decides how much to return}
}

\newglossaryentry{trustor}{
    name={Trustor},
    description={Player A in the Trust Game who decides how much money to send to the Trustee (also referred to as the first player)}
}

\newglossaryentry{trustee}{
    name={Trustee},
    description={Player B in the Trust Game who receives the multiplied amount and decides how much to return to the Trustor (also referred to as the second player)}
}

\newglossaryentry{archetype}{
    name={archetype},
    description={Standardized behavioral pattern or strategy used to define how agents act in specific scenarios, particularly for Trustee behavior in Trust Games}
}

% Methodology and evaluation
\newglossaryentry{belief-elicitation}{
    name={belief elicitation},
    description={The process of prompting language models to explicitly state their expectations or beliefs about behavioral patterns or outcomes}
}

\newglossaryentry{aggregate-level}{
    name={aggregate-level},
    description={Evaluation regime that examines statistical correlations between persona traits and behavioral outcomes across an entire population of agents}
}

\newglossaryentry{instance-level}{
    name={instance-level},
    description={Evaluation regime that assesses individual agents' ability to forecast their own specific actions in well-defined scenarios}
}

\newglossaryentry{effect-size}{
    name={effect size},
    description={Statistical measure of the strength of the relationship between variables, quantified using eta-squared ($\eta^2$) from one-way ANOVA}
}

\newglossaryentry{self-conditioning}{
    name={self-conditioning},
    description={Experimental condition where agents receive their own elicited beliefs as context to reinforce consistency between stated beliefs and behaviors}
}

% Technical terms and metrics
\newacronym{mae}{MAE}{Mean Absolute Error}

\newglossaryentry{spearman-correlation}{
    name={Spearman correlation},
    description={Non-parametric measure of rank correlation used to evaluate the consistency of trait-behavior relationship orderings},
    symbol={$\rho$}
}

\newglossaryentry{llm}{
    name={large language model},
    description={Neural network architecture trained on vast text corpora to generate human-like text responses},
    plural={large language models},
    symbol={LLM}
}

% Data and personas
\newglossaryentry{persona}{
    name={persona},
    description={Synthetic individual characterized by specific demographic, personality, and contextual attributes used in role-playing simulations}
}

\newglossaryentry{genagents}{
    name={GenAgents},
    description={Dataset of synthetic personas with demographic and social traits, augmented in this work with Big Five personality dimensions}
}

\newglossaryentry{big-five}{
    name={Big Five},
    description={Five-factor model of personality including Openness, Conscientiousness, Extraversion, Agreeableness, and Neuroticism}
}

% Relationships and patterns
\newglossaryentry{trait-behavior}{
    name={trait-behavior relationship},
    description={Statistical association between individual characteristics (demographic, personality) and behavioral outcomes in experimental settings}
}

\newglossaryentry{simulation-outcome}{
    name={simulation outcome},
    description={Behavioral results produced when language models role-play specific personas in experimental scenarios}
}

% Experimental conditions
\newglossaryentry{conditioning}{
    name={conditioning},
    description={Experimental manipulation where agents receive specific belief statements or instructions to guide their behavioral responses}
}

\newglossaryentry{perturbation}{
    name={perturbation},
    description={Systematic modification of elicited beliefs to test whether agents can be guided by researcher-imposed theoretical assumptions}
}